\section{Methods and Data}

\subsection{Patient and cell line RNA-seq data collection}

Raw RNA-seq FASTQ files from 158 human cancer cell lines were obtained from the Leibniz-Institut DSMZ – Deutsche Sammlung von Mikroorganismen und Zellkulturen GmbH (German Collection of Microorganisms and Cell Cultures). The cell lines represented multiple cancer types, including breast cancer, retinoblastoma, neuroblastoma, and haematological malignancies. RNA-seq count data for patient tumour samples were obtained from The Cancer Genome Atlas (TCGA) via the Genomic Data Commons (GDC) portal \cite{Geistlinger_Ramos_curatedTCGADataRef_2025}. Raw RNA-seq FASTQ files for additional patient tumour samples were obtained from the Gene Expression Omnibus (GEO).

\begin{table}[H]
\centering
\caption{Overview of RNA-seq datasets analysed in this study.}
\label{tab:dataset_summary}
\begin{tabular}{lcc}
\toprule
Cancer type & DSMZ cell lines & Patient tumour samples \\
\midrule
Breast cancer & 29 & 1231 \\
Neuroblastoma & 18 & 282 \\
Retinoblastoma & 11 & 66 \\
Haematological malignancies & 109 & 217 \\
\bottomrule
\end{tabular}
\end{table}


\begin{table}[H]
\centering
\caption{Public RNA-seq datasets used for patient tumour samples obtained from TCGA and GEO.}
\label{tab:patient_datasets}
\begin{tabular}{lll}
\toprule
Cancer type & Dataset accession & Samples \\
\midrule

Breast invasive carcinoma & TCGA-BRCA & 1231 \\

Neuroblastoma & GSE100148 & 4 \\
Neuroblastoma & GSE189367 & 64 \\
Neuroblastoma & GSE49711 & 5 \\
Neuroblastoma & SRP409177 & 47 \\
Neuroblastoma & TARGET-NBL & 162 \\

Retinoblastoma & GSE111168 & 6 \\
Retinoblastoma & GSE196420 & 18 \\
Retinoblastoma & GSE268136 & 42 \\

Burkitt lymphoma & GSE199413 & 8 \\
Chronic lymphocytic leukaemia & GSE237907 & 10 \\
Diffuse large B-cell lymphoma & TCGA-DLBC & 48 \\
Acute myeloid leukaemia & TCGA-LAML & 151 \\
Myeloproliferative neoplasms & GSE228758 & 6 \\

\bottomrule
\end{tabular}
\end{table}

\subsection{RNA-seq data cleaning and processing}

\subsubsection{Preprocessing steps and quality control}

Raw paired-end RNA-seq FASTQ files from DSMZ cell lines and patient tumour samples were processed using a workflow analogous to the GDC mRNA-seq pipeline \citep{GDCpipeline}.  Sequencing quality assessment was performed using FastQC (v0.11.9). Reads were aligned to the human reference genome (GRCh38/hg38) using the STAR aligner (v2.7.10b) \citep{Dobin2013} in two-pass mode. The resulting alignments were coordinate-sorted and indexed with SAMtools (v1.17) \citep{Danecek2021}. Library strandedness was inferred for each sample using RSeQC (v5.0.1) \citep{Wang2012} based on the aligned BAM files. Transcriptome-wide read coverage was computed using a custom Python script implemented with the \texttt{pysam} library (v0.21.0), generating per-sample coverage summaries. Gene annotations were obtained from GENCODE~v44 (GRCh38; Ensembl release~110). Alignment-level quality metrics were collected for each BAM file using Picard tools (v3.0+) \citep{Picard2019}. Project-level quality control reports were generated with MultiQC (v1.14) \citep{Ewels2016}, which aggregated FastQC results, STAR alignment statistics, and Picard metrics into unified interactive HTML reports.

\subsubsection{Tumour purity estimation}

Stromal and immune cells constitute the major non-tumour components of patient tumour samples. To estimate tumour purity, the ESTIMATE algorithm \cite{Yoshihara2013} was applied to infer stromal and immune cell infiltration from gene expression data.

To mitigate the confounding influence of non-tumour cellular admixture (e.g., immune and stromal cells), only patient tumour samples with a tumour purity score $\geq 0.7$ were retained for further analysis. This threshold was chosen to ensure that the analysed transcriptomic profiles predominantly reflect tumour-intrinsic gene expression patterns, which are more comparable to cancer cell lines that consist almost exclusively of malignant cells. Notably, tumour purity thresholds around 0.7 were evaluated during the validation of the ESTIMATE algorithm when comparing predicted purity with ABSOLUTE-derived tumour purity estimates \cite{Yoshihara2013}.

Pearson correlation coefficients were computed between tumour purity scores and the corresponding immune and stromal scores within each data source, and the results were visualised accordingly.
